\begin{abstract}
	
	\noindent
	My work is HUGE, I would need several days to explain it all, and even it has a lot of open branches and probably THOUSANDS of consequences. I designed an experiment, to prove international community it deserves to be studied and receive A LOT OF RESOURCES. Just in 5 hours. It is just a particular case, applied, of a more general technic. It has four phases.
	\\
	
	\noindent
	PHASE 01: A strange and very weird mathematical phenomenom, that is a system of relations very similar to an injective relation, by infinite parts, BUT IT IS NOT an injective relation, exactly.
	\\
	
	\noindent
	PHASE 02: DI. Inverse Diagonalization. (based in PHASE 01)
	\\
	
	\noindent
	PHASE 03: TPI. Unlimited Transfer of Pairs. (based in PHASE 01)
	\\
	
	\noindent
	PHASE 04: All participants, experts them all, will swear by the glory of their mother, that TPI or DI is a perfect builded counterexample of Cantor's Theorem. Trying to deny TPI or DI, they will GUARANTEE, exactly, the property that makes THE OTHER counterexample, a perfect injectivity practical equivalent:
	$$ r: \mathcal{P}(\mathbb{N}) \longrightarrow \mathbb{N}  $$
	\\\\
	
	\noindent
	My partner, mathematician, decided I was worthy just seeing PHASE 01. But I saw that other mathematicians, don't saw it as clear as him. So I tried hard for years...  and after many tries, many branches, I ended finding this strange phenomenom, where experts of all type and level, SWEAR, one of my two branches is correct. I created them to give PHASE 01 CLEAR MEANING.
	\\\\
	
	\noindent
	A couple of days ago, I found someone, an anonymous account in twitter, @AskPerplexity, that after a large talking... THEY GOT THE POINT OF THE MEANING of the system of relations. What it means related to Cantor's Theorem: it should not exists, if the theorem is right. And he asked me, JUST, for details of PHASE 01, very politely. I don't need PHASE 02, PHASE 03, or PHASE 04 with them !!. Even knowing, one step is ADMITING we are not using a proper injectivity. That we can challenge the definition of cardinality, with alternative tools to compare infinite cardinals.
	\\\\
	
	\noindent
	A perfect situation... the system of relations was checked by a lot of people before. It is very well constructed, and it does EXACTLY what I said them it can do. Other people never saw its meaning, or they did not have the courage to say it in public...
	\\\\
	
	\noindent
	This document contains the details, in my not academic language, about HOW TO BUILD that system of relations. I put another part, because we have an initial context, for those that could be a little lost, and to fix my bad habit of trying to explain too much. Like in this abstract :D.

\end{abstract}